% chapters/literature_review.tex
\chapter{Background and Related Work}
\label{ch:literature}

\section{Overview}
\label{sec:lit-overview}

This framework draws on three interconnected areas of research: the limitations of
LLM-generated causal explanations, human-centred explainability, and structured
evaluation methodology through multi-criteria decision making. This chapter reviews
each area and identifies the gaps that motivate the current work.

\section{Causal Reasoning in Large Language Models}
\label{sec:lit-causal}

Research on causal reasoning in LLMs is still developing, and current findings
indicate that models frequently conflate causation with correlation
\citep{jin2024llm, wu2024causality}. Large-scale pretraining enables models to
generate plausible causal explanations, but the reliability of these outputs for
academic and decision-support contexts remains uncertain \citep{kiciman2024causal}.

Studies have reported that LLMs can describe cause and effect at a surface level,
but their reasoning may shift across prompts or contradict known causal structures
\citep{miliani2026explica}. This variability makes evaluation important, particularly
when explanations are intended for decision support or for building student
understanding. The problem of evaluating causal explanation quality in LLMs is
therefore timely and practically relevant \citep{lin2024criticbench}.

\section{Human-Centred Explainability}
\label{sec:lit-xai}

Explainable AI (XAI) research has explored how people interpret explanations and
what makes them useful \citep{ehsan2020human, boydgraber2022human}. Prior work
identifies several key elements, including clarity, completeness, and support for
correct mental models. Other studies focus on trust, user satisfaction, and the
cognitive effort required to process explanations \citep{hoffman2019metrics}.

Research on how students and users engage with AI-generated content stresses that
explanations are effective only when they support understanding, trust, and sound
judgment \citep{ehsan2021expanding}. Surface fluency can mask weak or incomplete
reasoning, and this gap between fluency and quality is particularly consequential in
educational contexts where students may misplace trust in coherent-sounding but
causally flawed outputs.

Most current methods for human-centred evaluation remain qualitative or rely on
isolated criteria \citep{boydgraber2022human}. Comparisons across models are
therefore often inconsistent and difficult to generalise. The evaluation dimensions
adopted in the present study follow the framework of \citet{hoffman2019metrics},
specifically the four criteria of understandability, completeness, usefulness, and
trustworthiness.

\section{Structured Evaluation Methodology}
\label{sec:lit-mcdm}

Several approaches attempt to measure explanation quality, yet few target causal
explanations specifically, and fewer still combine multiple human-centred criteria
into one structured assessment \citep{wang2025llm}. Technical metrics tend to assess
model accuracy or alignment with ground truth, while user studies often depend on
open-ended feedback or simple ratings. These methods are helpful in isolation but do
not combine criteria in a principled way.

Multi-Criteria Decision Making (MCDM) offers tools for structured evaluation across
several fields, including operations research, engineering, and management science
\citep{mardani2015multiple}. These methods support the comparison of alternatives
against multiple criteria while incorporating expert or user judgment in a
transparent and reproducible way. Despite their wide use in assessment tasks, MCDM
methods have rarely been applied to the evaluation of LLM explanations.

The Analytic Hierarchy Process (AHP) \citep{saaty1987ahp} offers a principled
approach to deriving criterion weights from human judgment. It explicitly handles
hierarchical relationships between criteria and checks the logical consistency of
the provided judgments through the Consistency Ratio (CR). \citet{diaz2025ahp}
demonstrate the applicability of AHP to explainable AI frameworks in business
decision settings.

Interval-based reliability measures treat rating variability as genuine uncertainty
rather than noise \citep{hayes2007answering}. This perspective is consistent with
recent annotation research showing that reviewer disagreement carries meaningful
signal about the difficulty and ambiguity of the items being rated
\citep{beck2024order, zapf2016measuring}. Targeted sampling strategies
\citep{feng2025sample} further show that concentrating annotation effort on the most
discriminative cases improves evaluation efficiency without reducing coverage.

\section{Gap and Positioning}
\label{sec:lit-gap}

The present work addresses a clear gap: no existing framework combines
student-centred evaluation criteria, adaptive discriminative question selection, and
interval-based aggregation of reviewer disagreement into a single, reproducible
pipeline for evaluating LLM causal explanations. By drawing on AHP for criterion
weighting, semantic and structural dispersion for question selection, and IQR
intervals for aggregation, this study offers a more principled and transparent
comparison than prior work achieves with scalar metrics alone.
